% answers_short.tex —— 2006 年·四川卷·文科 · 简短答案速查
% 由 tools/make_short_answers.py 自 answers.tex 的 \gksolhead/\gkanswer 提取，
% 并与 2006/2006_数学真题_四川文_答案要点.md 逐题比对后生成，勿手改（改源 MD 后重跑）。

\begin{gknotice}
\textbf{使用说明}：本节只给最终答案便于快速核对，完整推导见解析版。答案为编写者独立推导，\textbf{非官方标准答案}。
\end{gknotice}

\gkpart{选择题}
\begin{center}\begin{tabular}{ccccccccccccc}
\toprule
题号 & 1 & 2 & 3 & 4 & 5 & 6 & 7 & 8 & 9 & 10 & 11 & 12 \\
\midrule
答案 & C & A & D & A & B & D & B & C & D & B & A & C \\
\bottomrule
\end{tabular}\end{center}

\gkpart{填空题}
\begin{itemize}[leftmargin=2.6em]
  \item[\textbf{13.}] $-960$
  \item[\textbf{14.}] $-6$
  \item[\textbf{15.}] $35$
  \item[\textbf{16.}] ①②④
\end{itemize}

\gkpart{解答题}
\begin{itemize}[leftmargin=2.6em]
  \item[\textbf{17.}] (1) $a_n=3^{n-1}$；(2) $T_n=n(n+2)$
  \item[\textbf{18.}] (1) $A=60^\circ$；(2) $\tan C=\dfrac{8+5\sqrt3}{11}$
  \item[\textbf{19.}] (1) $0.902$；(2) $0.254$
  \item[\textbf{20.}] (1) 见解析版；(2) $\arccos\dfrac2{\sqrt{21}}$
  \item[\textbf{21.}] (1) $x\in\left(-\dfrac23,1\right)$；(2) $|m|<\sqrt[3]{2}$
  \item[\textbf{22.}] (1) $-\sqrt2<k<-1$；(2) $m=4,\ S=\sqrt3$
\end{itemize}
